\chapter{Demonstration Chapter: Linear Systems}

\begin{learningobjectives}
By the end of this chapter, the reader should be able to:
\begin{itemize}
  \item identify a linear system and its coefficient matrix;
  \item state and use the invertibility criterion for uniqueness;
  \item distinguish a proof-based result from a computational procedure;
  \item recognise common algebraic errors in row reduction.
\end{itemize}
\end{learningobjectives}

\section{Core definitions}

\begin{definition}[Linear system]\label{def:linear-system}
A linear system of $m$ equations in $n$ unknowns is an equation of the form
\[
  A\bm{x}=\bm{b},
\]
where $A\in\R^{m\times n}$ is the coefficient matrix,
$\bm{x}\in\R^n$ is the vector of unknowns, and $\bm{b}\in\R^m$ is the data vector.
\end{definition}

\begin{remark}
Definitions should introduce notation before it is used. If a symbol is local to one
section, define it locally rather than adding it to the global notation list.
\end{remark}

\begin{keypoint}
\CoreTag\quad The equation $A\bm{x}=\bm{b}$ separates the fixed linear map $A$
from the input $\bm{b}$. This viewpoint is usually more reusable than treating each
scalar equation independently.
\end{keypoint}

\section{A representative theorem}

\begin{theorem}[Invertibility and uniqueness]\label{thm:invertibility-uniqueness}
Let $A\in\R^{n\times n}$. The following statements are equivalent:
\begin{enumerate}
  \item $A$ is invertible;
  \item for every $\bm{b}\in\R^n$, the system $A\bm{x}=\bm{b}$ has a unique solution;
  \item the homogeneous system $A\bm{x}=\bm{0}$ has only the trivial solution.
\end{enumerate}
\end{theorem}

\begin{proof}
Assume first that $A$ is invertible. Then for every $\bm{b}\in\R^n$,
\[
  \bm{x}=A^{-1}\bm{b}
\]
is a solution, and multiplying any solution by $A^{-1}$ shows that it is unique.
Hence (1) implies (2).

Statement (2) immediately implies (3) by taking $\bm{b}=\bm{0}$.

Finally, assume (3). Then the linear map $T(\bm{x})=A\bm{x}$ has trivial kernel, so
$T$ is injective. Since the domain and codomain both have dimension $n$, injectivity
implies surjectivity. Therefore $T$ is bijective and $A$ is invertible. Thus (3)
implies (1).
\end{proof}

\begin{examnote}
\ExamTag\quad In a proof question, write the logical direction being proved and name
the finite-dimensional result used in the last step. In a calculation question, it is
enough to use row reduction or a determinant criterion if the method is justified.
\end{examnote}

\section{Worked example}

\begin{example}[Solving a two-dimensional system]\label{ex:two-by-two}
Solve
\[
  \begin{pmatrix}
    2 & 1\\
    1 & -1
  \end{pmatrix}
  \begin{pmatrix}x\\y\end{pmatrix}
  =
  \begin{pmatrix}5\\1\end{pmatrix}.
\]
\end{example}

\begin{coursesolution}
The scalar equations are
\[
  2x+y=5,
  \qquad
  x-y=1.
\]
Adding the equations gives $3x=6$, hence $x=2$. Substitution into $x-y=1$ gives
$y=1$. Therefore
\[
  \bm{x}=\begin{pmatrix}2\\1\end{pmatrix}.
\]
A direct substitution verifies the result.
\end{coursesolution}

\begin{commonerror}
\CheckTag\quad Row operations must be applied to the entire augmented row, including
the right-hand side. A correct coefficient reduction paired with an unchanged data
column produces an invalid system.
\end{commonerror}

\section{Chapter summary}

\begin{itemize}
  \item A linear system is written compactly as $A\bm{x}=\bm{b}$.
  \item For square matrices, invertibility is equivalent to unique solvability for every
        right-hand side.
  \item Proofs should identify the structural reason for a conclusion; computations
        should record enough intermediate work to be checked.
\end{itemize}
